\section*{Ejemplo Coordenadas Cilíndricas - Ejercicio 25} 

\textcolor{red}{Ver comentarios sobre formato ejercicio 30.  Usar los mismos paquetes asi les es mas facil de unificar luego. }

\textbf{Enunciado:} Calcular el volumen del cilindro $x^2 + y^2 = 4$, bajo el paraboloide $2z + x^2 + y^2 = 16$ y sobre el plano $z = 0$.

\subsection*{1. Análisis de la región de integración $W$}

Definamos la región $W \subseteq \mathbb{R}^3$ como un conjunto elemental seg\'un la \textbf{Definición 2.5.7}.  


\textcolor{red}{Aca estaría bueno pegar la figura y el enlace al geogebra.} 


La proyección de este sólido sobre el plano $xy$ es el disco $$D_1 = \{ (x, y) \in \mathbb{R}^2 : x^2 + y^2 \leq 4 \}$$. 

Las superficies que limitan la altura $z$ son: \textcolor{red}{Ojo que las $\eta$ no son superficies, son funciones. Ver Definicion y escribir igual.  }
\begin{itemize}
    \item \textbf{Límite inferior ($\eta_{11}$):} El plano $z = 0$.
    \item \textbf{Límite superior ($\eta_{21}$):} El paraboloide $2z + x^2 + y^2 = 16$. Despejando $z$, obtenemos la función: 
    $$z = \frac{16 - (x^2 + y^2)}{2} = 8 - \frac{x^2 + y^2}{2}$$
\end{itemize}

Por lo tanto, el conjunto se expresa como:
$$W = \left\{ (x, y, z) \in \mathbb{R}^3 : (x, y) \in D_1 \wedge 0 \leq z \leq 8 - \frac{x^2 + y^2}{2} \right\}$$

\subsection*{2. Cambio de variables a coordenadas cilíndricas}

Para simplificar el cálculo, aplicamos el \textbf{Teorema 2.5.2} mediante la transformación $T(r, \theta, z) = (r \cos \theta, r \sin \theta, z)$. El valor absoluto del determinante jacobiano es $|J_T| = r$. \textcolor{red}{referencia: Ejemplo 2.4.4} 

\subsubsection*{Nuevos parámetros en $W^*$}
Definimos el nuevo conjunto de integración $W^*$ analizando las restricciones originales:

\begin{enumerate}
    \item \textbf{Rango de $\theta$:} Dado que el cilindro representa un disco completo en el plano $xy$, la región barre una vuelta completa sobre el eje $z$. Por lo tanto:
    $$0 \leq \theta \leq 2\pi$$

    \item \textbf{Rango de $r$:} La frontera lateral es $x^2 + y^2 = 4$. Al sustituir:
    $$(r \cos \theta)^2 + (r \sin \theta)^2 = 4 \implies r^2 = 4 \implies r = 2$$
    Como incluimos el interior, el radio varía de $0$ a $2$:
    $$0 \leq r \leq 2$$

    \item \textbf{Rango de $z$:} Sustituimos $x^2 + y^2 = r^2$ en la cota superior $\eta_{21}$:
    $$z = 8 - \frac{r^2}{2}$$
    Los límites para la variable $z$ en $W^*$ son:
    $$0 \leq z \leq 8 - \frac{r^2}{2}$$
\end{enumerate}

\subsection*{3. Cálculo del Volumen}

Utilizando la \textbf{Definición 2.5.6}, planteamos la integral triple:
$$\text{Vol}(W) = \iiint_W 1 \, dV = \int_{0}^{2\pi} \int_{0}^{2} \int_{0}^{8 - \frac{r^2}{2}} r \, dz \, dr \, d\theta$$  

\textcolor{red}{no separen en pasos asi queda mas resumido}

\textbf{Paso 1: Integración respecto a $z$}
$$\int_{0}^{8 - \frac{r^2}{2}} r \, dz = [rz]_{0}^{8 - \frac{r^2}{2}} = 8r - \frac{r^3}{2}$$

\textbf{Paso 2: Integración respecto a $r$}
$$\int_{0}^{2} \left( 8r - \frac{r^3}{2} \right) dr = \left[ 4r^2 - \frac{r^4}{8} \right]_{0}^{2} = (16 - 2) = 14$$

\textbf{Paso 3: Integración respecto a $\theta$}
$$\int_{0}^{2\pi} 14 \, d\theta = [14\theta]_{0}^{2\pi} = 28\pi$$

\textbf{Resultado final:}
$$\text{Vol}(W) = 28\pi$$

